% !TEX program = pdflatex
% Options for packages loaded elsewhere
\PassOptionsToPackage{unicode}{hyperref}
\PassOptionsToPackage{hyphens}{url}
\documentclass[
]{article}
\usepackage{xcolor}
\usepackage{amsmath,amssymb}
\usepackage[normalem]{ulem}
\newcommand{\cancel}[1]{\text{\sout{\ensuremath{#1}}}}
\setcounter{secnumdepth}{3} % number sections, subsections, subsubsections
\usepackage{iftex}
\ifPDFTeX
  \usepackage[T1]{fontenc}
  \usepackage[utf8]{inputenc}
  \usepackage{textcomp} % provide euro and other symbols
\else % if luatex or xetex
  \usepackage{unicode-math} % this also loads fontspec
  \defaultfontfeatures{Scale=MatchLowercase}
  \defaultfontfeatures[\rmfamily]{Ligatures=TeX,Scale=1}
\fi
\usepackage{lmodern}
\ifPDFTeX\else
  % xetex/luatex font selection
\fi
% Use upquote if available, for straight quotes in verbatim environments
\IfFileExists{upquote.sty}{\usepackage{upquote}}{}
% \IfFileExists{microtype.sty}{% use microtype if available
%   \usepackage[]{microtype}
%   \UseMicrotypeSet[protrusion]{basicmath} % disable protrusion for tt fonts
% }{}
\makeatletter
\@ifundefined{KOMAClassName}{% if non-KOMA class
  \IfFileExists{parskip.sty}{%
    \usepackage{parskip}
  }{% else
    \setlength{\parindent}{0pt}
    \setlength{\parskip}{6pt plus 2pt minus 1pt}}
}{% if KOMA class
  \KOMAoptions{parskip=half}}
\makeatother
\usepackage{graphicx}
\makeatletter
\newsavebox\pandoc@box
\newcommand*\pandocbounded[1]{% scales image to fit in text height/width
  \sbox\pandoc@box{#1}%
  \Gscale@div\@tempa{\textheight}{\dimexpr\ht\pandoc@box+\dp\pandoc@box\relax}%
  \Gscale@div\@tempb{\linewidth}{\wd\pandoc@box}%
  \ifdim\@tempb\p@<\@tempa\p@\let\@tempa\@tempb\fi% select the smaller of both
  \ifdim\@tempa\p@<\p@\scalebox{\@tempa}{\usebox\pandoc@box}%
  \else\usebox{\pandoc@box}%
  \fi%
}
% Set default figure placement to htbp
\def\fps@figure{htbp}
\makeatother
\setlength{\emergencystretch}{3em} % prevent overfull lines
\providecommand{\tightlist}{%
  \setlength{\itemsep}{0pt}\setlength{\parskip}{0pt}}
\usepackage{bookmark}
\IfFileExists{xurl.sty}{\usepackage{xurl}}{} % add URL line breaks if available
\urlstyle{same}
\hypersetup{
  hidelinks,
  pdfcreator={LaTeX via pandoc}}
\let\oldhyperref\hyperref
\renewcommand{\hyperref}[2][]{\oldhyperref[#1]{\underline{#2}}}

\author{}
\date{}

\usepackage[margin=1in]{geometry}
\usepackage{multicol}
\usepackage{fancyhdr}
\pagestyle{fancy}
\fancyhf{}
\fancyhead[L]{\textbf{ECE 128 - HW5}}
\fancyhead[R]{Chaidhat Chaimongkol}
\fancyfoot[C]{\thepage}
\renewcommand{\headrulewidth}{0.4pt}
\providecommand{\ket}[1]{\left\lvert #1 \right\rangle}
\providecommand{\bra}[1]{\left\langle #1 \right\rvert}
\providecommand{\braket}[1]{\left\langle #1 \right\rangle}
\providecommand{\Rarr}{\Rightarrow}
\providecommand{\Larr}{\Leftarrow}
\providecommand{\lrarr}{\Leftrightarrow}
\providecommand{\rarr}{\rightarrow}
\providecommand{\larr}{\leftarrow}
\providecommand{\lrar}{\leftrightarrow}
\providecommand{\Z}{\mathbb{Z}}
\providecommand{\R}{\mathbb{R}}
\providecommand{\N}{\mathbb{N}}
\providecommand{\C}{\mathbb{C}}
\providecommand{\Q}{\mathbb{Q}}
\providecommand{\infin}{\infty}

\date{\today}


\begin{document}
\section*{Problem 1}
Let

\[
	\ket{\psi}=a\ket{0}_A+b\ket{1}_A,\qquad \ket{\phi}=c\ket{0}_B+d\ket{1}_B,
\]
\[
	|a|^2+|b|^2=1\quad |c|^2+|d|^2=1
\]
Note that
\[
	\ket{\psi,\phi}=\ket{\psi}\otimes\ket{\phi}=ac\ket{00}+ad\ket{01}+bc\ket{10}+bd\ket{11}
\]
Starting from this
\[
	\quad \braket{\Phi_+^{(AB)}|\psi,\phi}=\frac{1}{\sqrt{2}}\left(\bra{00}+\bra{11}\right)\left(ac\ket{00}+ad\ket{01}+bc\ket{10}+bd\ket{11}\right)
\]
\[
	=\frac{1}{\sqrt{2}}\Big[
	ac\braket{00|00}+\cancel{ad\braket{00|01}}+\cancel{bc\braket{00|10}}+\cancel{bd\braket{00|11}}
\]
\[
	+\,\cancel{ac\braket{11|00}}+\cancel{ad\braket{11|01}}+\cancel{bc\braket{11|10}}+bd\braket{11|11}\Big]
\]
\[
	=\frac{1}{\sqrt{2}}\left(ac+bd\right)
\]

Applying Cauchy-Schwarz
\[
	\left|ac+bd\right|^2 \leq \left(|a|^2+|b|^2\right)\left(|c|^2+|d|^2\right)=1\cdot 1 = 1
\]
Therefore $\left|ac+bd\right|\leq 1$, and hence
\[
	\left|\braket{\Phi_+^{(AB)}|\psi,\phi}\right|=\frac{|ac+bd|}{\sqrt{2}}\leq\frac{1}{\sqrt{2}}<1
\]

So $\left|\braket{\Phi_+^{(AB)}|\psi,\phi}\right|$ never can be equal to 1 so it is not a product state, it is an entangled state.

\section*{Problem 2}

\subsection*{Problem 2a}
$A=A$ always, so the probability not being itself is zero.
\[
	d(A,A)=\Pr(A\neq A)=0
\]
Probability of two things are commutative so $Pr(a,b)=Pr(b,a)$ so
\[
	d(A,B)=d(B,A)
\]
We got to think of it using set theory so if $A\neq C$ then either $A\neq B$ or $B\neq C$, so $\{A\neq C\}\subseteq\{A\neq B\}\cup\{B\neq C\}$. By the union bound
\[
	d(A,C)\leq d(A,B)+d(B,C)
\]

\subsection*{Problem 2b}
$A,B\in\{-1,+1\}$ gives $AB=+1$ if $A=B$, $AB=-1$ if $A\neq B$, so
\[
	\mathbb{E}[AB]=\Pr(A=B)-\Pr(A\neq B)=1-2\,d(A,B)
\]
\[
	\Rarr\quad d(A,B)=\frac{1-\mathbb{E}[AB]}{2}
\]

\subsection*{Problem 2c}
From 2a
\[
	d(A_1,A_2)\leq d(A_1,B_1)+d(B_1,A_2)
\]
\[
	d(A_1,B_2)\leq d(A_1,A_2)+d(A_2,B_2)
\]
Substitute
\[
	d(A_1,B_2)\leq d(A_1,B_1)+d(B_1,A_2)+d(A_2,B_2)
\]
This must be true.
So we know
\[
	A_1\;\rarr^{\;d(A_1,B_1)\;}\;B_1\;\rarr^{\;d(B_1,A_2)\;}\;A_2\;\rarr^{\;d(A_2,B_2)\;}\;B_2
\]
That means $A_1B_2$ is the diagonal.
Substituting $d=(1-\mathbb{E}[XY])/2$ from Problem 2b recovers
\[
	\mathbb{E}[A_1B_2]-\mathbb{E}[A_1B_1]-\mathbb{E}[B_1A_2]-\mathbb{E}[A_2B_2]\leq 2
\]

\subsection*{Problem 2d}
Let $\ket{\Psi_-}$,
\[
	\mathbb{E}[(\vec{\sigma}\cdot\hat{a})\otimes(\vec{\sigma}\cdot\hat{b})]=-\cos\theta_{ab}\quad d(A,B)=\frac{1+\cos\theta_{ab}}{2}
\]
Take $\hat{a}_1,\hat{b}_1,\hat{a}_2,\hat{b}_2$ on the same plane at angles
\[
	\alpha_{a_1}=0\quad \alpha_{b_1}=\tfrac{3\pi}{4}\quad \alpha_{a_2}=\tfrac{3\pi}{2}\quad \alpha_{b_2}=\tfrac{\pi}{4}
\]
The angle between two coplanar unit vectors is $\theta=\min(|\Delta\alpha|,\,2\pi-|\Delta\alpha|)$, so
\[
	\theta_{a_1b_1}=|0-\tfrac{3\pi}{4}|=\tfrac{3\pi}{4}
\]
\[
	\theta_{b_1a_2}=|\tfrac{3\pi}{4}-\tfrac{3\pi}{2}|=|\tfrac{3\pi}{4}-\tfrac{6\pi}{4}|=\tfrac{3\pi}{4}
\]
\[
	\theta_{a_2b_2}=2\pi-|\tfrac{3\pi}{2}-\tfrac{\pi}{4}|=2\pi-|\tfrac{6\pi}{4}-\tfrac{\pi}{4}|=2\pi-\tfrac{5\pi}{4}=\tfrac{3\pi}{4}
\]
\[
	\theta_{a_1b_2}=|0-\tfrac{\pi}{4}|=\tfrac{\pi}{4}
\]
Plug into $d(A,B)=\frac{1+\cos\theta_{ab}}{2}$, using $\cos\tfrac{3\pi}{4}=-\tfrac{1}{\sqrt{2}}$ and $\cos\tfrac{\pi}{4}=\tfrac{1}{\sqrt{2}}$
\[
	d(A_1,B_1)=\frac{1+\cos(3\pi/4)}{2}=\frac{1+(-1/\sqrt{2})}{2}=\frac{1-1/\sqrt{2}}{2}
\]
\[
	d(B_1,A_2)=\frac{1+\cos(3\pi/4)}{2}=\frac{1-1/\sqrt{2}}{2}
\]
\[
	d(A_2,B_2)=\frac{1+\cos(3\pi/4)}{2}=\frac{1-1/\sqrt{2}}{2}
\]
\[
	d(A_1,B_2)=\frac{1+\cos(\pi/4)}{2}=\frac{1+1/\sqrt{2}}{2}
\]
From 2c,
\[
	d(A_1,B_2)\leq d(A_1,B_1)+d(B_1,A_2)+d(A_2,B_2)
\]
\[
	\;0.85\approx\frac{1+1/\sqrt{2}}{2}\nleq
	\frac{3-3/\sqrt{2}}{2}\approx 0.44\;
\]
Hence contradiction, so it is not possible.

\section*{Problem 3}

\subsection*{Problem 3a}
$\sigma_Z$ has eigenvectors $+1$ for $\ket{0}$ and $-1$ for $\ket{1}$ so
\[
	\ket{\text{GHZ}}=\frac{1}{\sqrt{2}}(\ket{0}_A\ket{00}_{BC}-\ket{1}_A\ket{11}_{BC})
\]
Alice measures $+1$ $\Rarr$ project onto $\ket{0}_A$
\[
	\ket{\text{GHZ}}\rarr\ket{0}_A\ket{00}_{BC}=\ket{0}_A\otimes\ket{0}_B\otimes\ket{0}_C
\]
Alice measures $-1$ $\Rarr$ project onto $\ket{1}_A$
\[
	\ket{\text{GHZ}}\rarr-\ket{1}_A\ket{11}_{BC}=-\ket{1}_A\otimes\ket{1}_B\otimes\ket{1}_C
\]
Either way the BC subsystem is a product of single-qubit states.

\subsection*{Problem 3b}
$\sigma_X$ has eigenvectors $\ket{\pm}=\frac{1}{\sqrt{2}}(\ket{0}\pm\ket{1})$, so
\[
	\ket{0}=\frac{\ket{+}+\ket{-}}{\sqrt{2}}\quad \ket{1}=\frac{\ket{+}-\ket{-}}{\sqrt{2}}
\]
Sub into $\ket{\text{GHZ}}$
\[
	\ket{\text{GHZ}}=\frac{1}{\sqrt{2}}\left[\frac{\ket{+}_A+\ket{-}_A}{\sqrt{2}}\ket{00}_{BC}-\frac{\ket{+}_A-\ket{-}_A}{\sqrt{2}}\ket{11}_{BC}\right]
\]
\[
	=\frac{1}{2}\Big[\ket{+}_A(\ket{00}-\ket{11})_{BC}+\ket{-}_A(\ket{00}+\ket{11})_{BC}\Big]
\]
Alice measures $+1$ $\Rarr$ BC $=\frac{1}{\sqrt{2}}(\ket{00}-\ket{11})=\ket{\Phi_-^{(BC)}}$.

Alice measures $-1$ $\Rarr$ BC $=\frac{1}{\sqrt{2}}(\ket{00}+\ket{11})=\ket{\Phi_+^{(BC)}}$.

Both are Bell states. By Schwartz like in 1),
\[
	\left|\braket{\Phi_\pm^{(BC)}|\psi,\phi}\right|=\frac{|ac\pm bd|}{\sqrt{2}}\leq\frac{1}{\sqrt{2}}<1
\]
But if BC were $\ket{\psi}\otimes\ket{\phi}$ the overlap would equal $1$.

Contradiction thus BC is entangled.

\section*{Problem 4}

\subsection*{Problem 4a}
For arbitrary $\ket{\psi},\ket{\phi},\ket{\psi'},\ket{\phi'}$
\[
	\bra{\psi',\phi'}(A\otimes B)^+\ket{\psi,\phi}=\braket{\psi,\phi|A\otimes B|\psi',\phi'}^*
\]
\[
	=\big(\bra{\psi}A\ket{\psi'}\bra{\phi}B\ket{\phi'}\big)^*=\bra{\psi}A\ket{\psi'}^*\bra{\phi}B\ket{\phi'}^*
\]
\[
	=\bra{\psi'}A^+\ket{\psi}\bra{\phi'}B^+\ket{\phi}=\bra{\psi',\phi'}A^+\otimes B^+\ket{\psi,\phi}
\]
Matrix elements agree on all states, so
\[
	(A\otimes B)^+=A^+\otimes B^+
\]

\subsection*{Problem 4b}
$A^+=A$, $B^+=B$, so by 4a
\[
	(A\otimes B)^+=A^+\otimes B^+=A\otimes B
\]

\subsection*{Problem 4c}
Use $(X\otimes Y)(X'\otimes Y')=(XX')\otimes(YY')$ and 4a throughout

Positive: $A\triangleq M^+ M$, $B\triangleq N^+ N$, so
\[
	A\otimes B=(M^+ M)\otimes(N^+ N)=(M\otimes N)^+(M\otimes N)\geq 0
\]

Unitary: $A^+ A=I$, $B^+ B=I$, so
\[
	(A\otimes B)^+(A\otimes B)=(A^+\otimes B^+)(A\otimes B)=(A^+ A)\otimes(B^+ B)=I\otimes I=I
\]

Normal: $A^+ A=AA^+$, $B^+ B=BB^+$, so
\[
	(A\otimes B)^+(A\otimes B)=(A^+ A)\otimes(B^+ B)=(AA^+)\otimes(BB^+)=(A\otimes B)(A\otimes B)^+
\]

\subsection*{Problem 4d}
No. Take $A=B=iI$. Neither is hermitian ($A^+=-iI\neq A$), but by 4a
\[
	(A\otimes B)^+=A^+\otimes B^+=(-iI)\otimes(-iI)=-I\otimes I=A\otimes B
\]
which is hermitian. More generally, anti-herm$\otimes$ anti-herm is herm
\[
	(A\otimes B)^+=(-A)\otimes(-B)=A\otimes B
\]

\section*{Problem 5}

\subsection*{Problem 5a}
$R$ is unitary, so preserves inner products
\[
	\braket{R0|R1}=\braket{0|1}=0
\]
Hence $R\ket{1}\perp R\ket{0}=\ket{n_+}$. In a 2D Hilbert space there is exactly one direction orthogonal to $\ket{n_+}$, namely $\ket{n_-}$ (eigenvectors are unique only up to a phase), so
\[
	R\ket{1}=e^{i\beta}\ket{n_-}
\]

\subsection*{Problem 5b}
Allow arbitrary phases $R\ket{0}=e^{i\alpha}\ket{n_+}$, $R\ket{1}=e^{i\beta}\ket{n_-}$. Apply $R\otimes R$ to $\ket{\Psi_-}$
\[
	(R\otimes R)\ket{\Psi_-}=\frac{1}{\sqrt{2}}\big[(R\ket{0})\otimes(R\ket{1})-(R\ket{1})\otimes(R\ket{0})\big]
\]
\[
	=\frac{e^{i(\alpha+\beta)}}{\sqrt{2}}\big[\ket{n_+}\otimes\ket{n_-}-\ket{n_-}\otimes\ket{n_+}\big]
\]
Expand $\ket{n_+}=c\ket{0}+s\ket{1}$ and $\ket{n_-}=-s^*\ket{0}+c\ket{1}$ with $c=\cos(\theta/2)$, $s=e^{i\phi}\sin(\theta/2)$, $c^2+|s|^2=1$
\[
	\ket{n_+}\otimes\ket{n_-}=-cs^*\ket{00}+c^2\ket{01}-|s|^2\ket{10}+cs\ket{11}
\]
\[
	\ket{n_-}\otimes\ket{n_+}=-s^*c\ket{00}-|s|^2\ket{01}+c^2\ket{10}+cs\ket{11}
\]
Subtract
\[
	\ket{n_+}\otimes\ket{n_-}-\ket{n_-}\otimes\ket{n_+}=(c^2+|s|^2)\ket{01}-(|s|^2+c^2)\ket{10}=\ket{01}-\ket{10}=\sqrt{2}\,\ket{\Psi_-}
\]
\[
	\Rarr\quad (R\otimes R)\ket{\Psi_-}=e^{i(\alpha+\beta)}\ket{\Psi_-}
\]
i.e.~$\ket{\Psi_-}$ up to a phase.

\subsection*{Problem 5c}
Let $R$ be the rotation taking $\hat{z}\to\hat{n}$. Then $R^+\ket{n_\pm}=\ket{0/1}$ (up to phase) and the eigenvalues match, so $R^+ S_n R=S_z$, i.e.
\[
	S_n=R\,S_z\,R^+
\]
Therefore
\[
	\bra{\Psi_-}S_n\otimes S_n\ket{\Psi_-}=\bra{\Psi_-}(R\,S_z\,R^+)\otimes(R\,S_z\,R^+)\ket{\Psi_-}
\]
\[
	=\bra{\Psi_-}(R\otimes R)(S_z\otimes S_z)(R^+\otimes R^+)\ket{\Psi_-}
\]
From (b), $(R^+\otimes R^+)\ket{\Psi_-}=e^{-i(\alpha+\beta)}\ket{\Psi_-}$ and $\bra{\Psi_-}(R\otimes R)=e^{i(\alpha+\beta)}\bra{\Psi_-}$. The phases cancel out and we get
\[
	\bra{\Psi_-}S_n\otimes S_n\ket{\Psi_-}=\bra{\Psi_-}S_z\otimes S_z\ket{\Psi_-}
\]
Compute the RHS using $S_z\ket{0}=\tfrac{\hbar}{2}\ket{0}$ and $S_z\ket{1}=-\tfrac{\hbar}{2}\ket{1}$
\[
	(S_z\otimes S_z)\ket{01}=-\tfrac{\hbar^2}{4}\ket{01}\quad (S_z\otimes S_z)\ket{10}=-\tfrac{\hbar^2}{4}\ket{10}
\]
\[
	(S_z\otimes S_z)\ket{\Psi_-}=\frac{1}{\sqrt{2}}\left[-\tfrac{\hbar^2}{4}\ket{01}+\tfrac{\hbar^2}{4}\ket{10}\right]=-\tfrac{\hbar^2}{4}\ket{\Psi_-}
\]
\[
	\Rarr\quad \bra{\Psi_-}S_z\otimes S_z\ket{\Psi_-}=-\tfrac{\hbar^2}{4}
\]
\[
	\frac{4}{\hbar^2}\bra{\Psi_-}S_n\otimes S_n\ket{\Psi_-}=\frac{4}{\hbar^2}\cdot\left(-\tfrac{\hbar^2}{4}\right)=-1
\]

\subsection*{Problem 5d}
DNE. If parallel measurements always got the same results, then $S_n\otimes S_n$ always returns $+\hbar^2/4$ (either both $+\hbar/2$ or both $-\hbar/2$), so
\[
	\bra{\psi}S_n\otimes S_n\ket{\psi}=+\tfrac{\hbar^2}{4}\quad\forall\hat{n}
\]
Sum over orthogonal axes $\hat{x},\hat{y},\hat{z}$
\[
	\bra{\psi}S_x\otimes S_x\ket{\psi}+\bra{\psi}S_y\otimes S_y\ket{\psi}+\bra{\psi}S_z\otimes S_z\ket{\psi}=\tfrac{3\hbar^2}{4}
\]
The LHS equals $\bra{\psi}\vec{S}_1\vec{S}_2\ket{\psi}$. Use the total spin $\vec{S}=\vec{S}_1+\vec{S}_2$
\[
	\vec{S}^2=\vec{S}_1^2+\vec{S}_2^2+2\,\vec{S}_1\vec{S}_2=\tfrac{3\hbar^2}{2}+2\,\vec{S}_1\vec{S}_2
\]
For two spin-$\tfrac{1}{2}$ particles, $\vec{S}^2$ has eigenvalues $0$ (singlet) or $2\hbar^2$ (triplet), so
\[
	\vec{S}_1\vec{S}_2\in\left\{-\tfrac{3\hbar^2}{4},\;+\tfrac{\hbar^2}{4}\right\}
\]
Max. eigenvalue is $\hbar^2/4$, so $\bra{\psi}\vec{S}_1\vec{S}_2\ket{\psi}\leq\hbar^2/4$ for any pure state (and $\mathrm{Tr}(\rho\,\vec{S}_1\vec{S}_2)\leq\hbar^2/4$ for any mixed state). But the anti-singlet would need $\tfrac{3\hbar^2}{4}>\tfrac{\hbar^2}{4}$.

Contradiction.

\end{document}
